\documentclass[10pt,twocolumn,letterpaper]{article}

\usepackage[utf8]{inputenc}
\usepackage{amsmath,amssymb,amsfonts}
\usepackage{graphicx}
\usepackage{booktabs}
\usepackage{hyperref}
\usepackage{microtype}
\usepackage{cite}
\usepackage[margin=0.75in]{geometry}

\title{\textbf{Kymora: High-Throughput Batch and Streaming Time-Series Feature Extraction with a Rust Core}}

\author{
  \textbf{Aamod Kumar} \\
  Independent Researcher \\
  \texttt{aamod@example.com}
}

\date{November 22, 2026}

\begin{document}

\maketitle

\begin{abstract}
Time-series feature extraction is an essential precursor to classification, clustering, anomaly detection, and regression in industrial IoT, biosignal monitoring, and quantitative finance. Established toolkits---notably \texttt{tsfresh}, \texttt{catch22}, and \texttt{TSFEL}---provide rich feature banks but incur steep computational overheads on large datasets due to per-series Python dispatch loops, repeated data allocation, and redundant intermediate transformations. We present \textbf{Kymora}, an open-source, high-throughput feature extraction library featuring an optimized Rust core exposed via PyO3. By borrowing contiguous NumPy memory buffers directly with zero intermediate allocations, releasing the Python Global Interpreter Lock (GIL), and scheduling work-stealing parallelism over series via Rayon, Kymora achieves a pooled headline throughput of 514,245 series/s (1.94~ms for 1,000 series $\times$ 500 steps on a 16-thread system). Under like-for-like equal-feature conditions, Kymora is 6.0$\times$ faster than Numba, 165.1$\times$ faster than vectorized NumPy, 1,089.1$\times$ faster than TSFEL, and 2,812.4$\times$ faster than tsfresh. Furthermore, across 42 canonical univariate UCR classification benchmarks, Kymora core33 matches catch22 predictive accuracy (79.96\% vs 79.83\% Random Forest accuracy; Holm-adjusted $p=0.5554$) while delivering a 117.2$\times$ extraction speedup. For sliding windows, Kymora provides an anchored streaming engine with amortized $O(1)$ per-sample updates. We report our full empirical loss ledger alongside wins to provide transparent engineering baselines.
\end{abstract}

\section{Introduction}
Modern sensor networks and telemetry systems ingest continuous high-dimensional time-series data. Extracting summary representations that capture central tendency, dispersion, autocorrelation, spectral density, and non-linear complexity transforms unstructured sequences into compact vectors for standard tabular learners~\cite{christ2018tsfresh,lubba2019catch22,barandas2020tsfel}.

Despite the utility of feature-based representation, practitioners face three operational bottlenecks:
\begin{enumerate}
    \item \textbf{FFI Dispatch Overhead:} Iterating over thousands of independent series in interpreted Python invokes tens of thousands of function calls and C foreign-function boundary crossings.
    \item \textbf{Redundant Transformations:} Independent feature functions repeatedly recompute intermediate representations (such as sample means, sorted values, difference series, or Fourier transforms), multiplying memory bandwidth demands.
    \item \textbf{Window Recomputation:} Streaming applications frequently re-evaluate full rolling windows from scratch ($O(W)$ per sample), creating computational latency spikes in real-time inference.
\end{enumerate}

Kymora resolves these bottlenecks through a fused Rust computation engine. It presents a clean Python API compatible with \texttt{scikit-learn}~\cite{pedregosa2011scikit} and Polars, while delegating heavy arithmetic to compiled, vectorizable Rust kernels.

\section{System Architecture}
\subsection{Direct Buffer Borrowing and Threading}
Kymora requires $C$-contiguous float64 or float32 arrays. PyO3 bindings inspect input buffers, validating strides and dimensionality in under 100~ns. Memory pointers are borrowed directly without copying data. The Python GIL is immediately released, allowing Rayon's work-stealing thread pool~\cite{rayon} to shard batches across worker threads.

To eliminate thread coordination overhead on small inputs, Kymora incorporates a sequential fast path: when total points $N \times L < 8,000$, execution runs inline on the caller thread. Single-series inference via \texttt{extract\_one} executes in 3.90~$\mu$s.

\subsection{Fused Kernel Passes}
For each series of length $L$, Kymora minimizes passes over memory:
\begin{itemize}
    \item \textbf{Pass 1 (Fused Accumulation):} Computes sum, sum-of-squares, minimum, maximum, and validates finiteness in a single cache-resident traversal.
    \item \textbf{Pass 2 (Centered Moments):} Computes sample variance, skewness, and kurtosis using centered differences.
    \item \textbf{Spectral & Nonlinear Analysis:} Computes real-to-complex FFT via \texttt{realfft}~\cite{frigo2005design}, deriving spectral centroids and band powers from a shared power spectrum. Permutation entropy~\cite{bandt2002permutation} and complexity-invariant distance~\cite{batista2014cid} are computed using reusable scratch arenas.
\end{itemize}

\subsection{Streaming Engine with Drift Guards}
For rolling windows of length $W$, \texttt{StreamingExtractor} implements shifted power sums, updating mean, variance, skewness, and kurtosis in amortized $O(1)$ time per pushed sample. To prevent catastrophic floating-point cancellation, an internal drift guard triggers an exact $O(W)$ re-anchor if accumulated numerical drift exceeds $0.25\sigma$. Benchmarks verify push latency at 0.20~$\mu$s p50 across window capacities from $W=64$ to $W=65,536$.

\section{Empirical Evaluation}
\subsection{Experimental Setup}
All timing benchmarks were conducted on an Intel Core i7-13620H system (10 cores [6P+4E], 16 threads, Windows 11) using Python 3.14.6 and NumPy 2.4.6 under an audited High-Performance scheme with inter-session cooldowns. Every reported metric traces to committed artifacts in \texttt{CLAIMS.md}.

\subsection{Like-for-Like Equal-Feature Extraction}
To eliminate discrepancies caused by differing feature counts, each library was constrained to definition-matched subsets, with output values parity-gated at relative error $\le 10^{-9}$ before timing.

\begin{table}[h]
\centering
\small
\caption{Equal-feature benchmark (1,000 series $\times$ 500 steps, 16 threads). Latencies are pooled medians.}
\label{tab:equal_feature}
\begin{tabular}{lrrrr}
\toprule
\textbf{Library} & \textbf{Feat.} & \textbf{Kymora (ms)} & \textbf{Target (ms)} & \textbf{Ratio} \\
\midrule
Numba baseline   & 33 & 2.24 & 13.34     & 6.0$\times$ \\
NumPy baseline   & 33 & 2.19 & 361.33    & 165.1$\times$ \\
TSFEL            & 13 & 1.97 & 2,145.73  & 1,089.1$\times$ \\
tsfresh          & 13 & 2.34 & 6,580.57  & 2,812.4$\times$ \\
\bottomrule
\end{tabular}
\end{table}

As shown in Table~\ref{tab:equal_feature}, Kymora maintains a 6.0$\times$ to 2,812.4$\times$ throughput advantage under identical mathematical specifications.

\subsection{Default Configuration Throughput}
In full default configurations (Kymora core33 vs.\ catch22 [22], TSFEL [156], and tsfresh [777]):
\begin{itemize}
    \item \textbf{Kymora core33:} 1.94~ms (514,245 series/s, 0.0589~$\mu$s/series-feature; isolated baseline 358,654 series/s).
    \item \textbf{catch22:} 185.50~ms (5,391 series/s), raw ratio 95.4$\times$ (143.1$\times$ per-feature).
    \item \textbf{TSFEL:} 2,509.53~ms (398 series/s), raw ratio 1,291$\times$ (273.0$\times$ per-feature).
    \item \textbf{tsfresh:} 23,540.62~ms (42 series/s), raw ratio 12,106$\times$ (514.1$\times$ per-feature).
\end{itemize}

\subsection{Downstream Classification on 42 UCR Datasets}
To assess whether high-speed features preserve discriminative fidelity, we evaluated classification performance across 42 canonical univariate UCR time-series datasets~\cite{dau2019ucr} using canonical train/test splits.

\begin{table}[h]
\centering
\small
\caption{Downstream UCR classification accuracy and extraction latency across 42 datasets (16 threads).}
\label{tab:ucr}
\begin{tabular}{lrrr}
\toprule
\textbf{Feature Set} & \textbf{RF Acc.} & \textbf{Ridge Acc.} & \textbf{Extr. Time} \\
\midrule
Kymora core33   & 79.96\% & 77.42\% & 1.26~ms \\
catch22         & 79.83\% & 75.75\% & 159.40~ms \\
TSFEL           & 81.00\% & 79.71\% & 942.67~ms \\
tsfresh         & 84.79\% & 82.39\% & 20,943.62~ms \\
\midrule
Kymora extended & 83.12\% & 80.39\% & 1.74~ms \\
Kymora full     & 82.79\% & 80.66\% & 2.17~ms \\
\bottomrule
\end{tabular}
\end{table}

Kymora core33 achieves catch22-class accuracy (79.96\% vs.\ 79.83\% Random Forest; head-to-head 23 wins, 3 ties, 16 losses; Holm-adjusted $p=0.5554$, indicating no statistically significant difference) while executing 117.2$\times$ faster. Wider feature banks (tsfresh and TSFEL) yield higher classification accuracy, but require up to 17,217.0$\times$ greater extraction time. Kymora's \texttt{extended} profile reaches 83.12\% accuracy in 1.74~ms.

\section{Limitations \& Loss Ledger}
We explicitly publish configurations where alternatives are preferable:
\begin{enumerate}
    \item \textbf{Micro-Batches:} On single-series micro-workloads ($1 \times 10$ samples), hand-tuned Numba achieves 3.00~$\mu$s versus Kymora's 5.90~$\mu$s, winning by approximately 2.0$\times$ due to zero FFI overhead.
    \item \textbf{Exploratory Screening:} When seeking hundreds of exploratory hypothesis-test features, tsfresh's 777-feature bank remains superior.
    \item \textbf{P/E Core Scaling:} On hybrid Intel architectures, Rayon thread scaling plateaus beyond physical Performance cores (8T at 2.22~ms vs.\ 16T at 2.30~ms).
\end{enumerate}

\section{Conclusion}
Kymora demonstrates that fusing time-series traversals within a compiled Rust core enables orders-of-magnitude extraction throughput gains while maintaining predictive accuracy on standard machine learning benchmarks.

\bibliographystyle{IEEEtran}
\bibliography{references}

\end{document}
